\documentclass[12pt]{article}

\usepackage[a4paper,margin=2.5cm]{geometry}
\usepackage{graphicx}
\usepackage{booktabs}
\usepackage{tcolorbox}
\usepackage{xcolor}
\usepackage{hyperref}
\usepackage{longtable}
\usepackage{enumitem}

\definecolor{primary}{HTML}{1F4E79}
\definecolor{softblue}{HTML}{E8F1FA}

\title{\textbf{\textcolor{primary}{Especificación de Requisitos de Software}}\\
\vspace{0.3cm}
JustBeat-DAW\\
\large 8-bit Digital Audio Workstation}
\author{Equipo de Desarrollo}
\date{\today}

\begin{document}

\maketitle
\tableofcontents
\newpage

%====================================================
\section{Introducción}

\subsection{Propósito}

Este documento define los requisitos funcionales y no funcionales del sistema \textbf{JustBeat-DAW}, una Digital Audio Workstation orientada a producción musical 8-bit (chiptune).

El documento servirá como base contractual para diseño, desarrollo, pruebas y evolución futura del sistema.

\subsection{Alcance}

JustBeat-DAW permitirá:

\begin{itemize}
\item Crear y gestionar proyectos musicales.
\item Generar sonido 8-bit mediante síntesis digital.
\item Utilizar entrada y grabación MIDI.
\item Soportar sistema extensible de plugins.
\item Exportar audio y MIDI.
\item Gestionar múltiples pistas y buses de audio.
\end{itemize}

\subsection{Definiciones}

\begin{itemize}
\item \textbf{DAW}: Digital Audio Workstation.
\item \textbf{Chiptune}: Música basada en síntesis de consolas retro.
\item \textbf{MIDI}: Musical Instrument Digital Interface.
\item \textbf{Plugin}: Módulo extensible que añade funcionalidad.
\item \textbf{Audio Graph}: Grafo de procesamiento de señal.
\end{itemize}

%====================================================
\section{Descripción General}

\subsection{Perspectiva del Producto}

Aplicación de escritorio modular con:

\begin{itemize}
\item Arquitectura limpia.
\item Motor de audio propio.
\item Soporte MIDI.
\item Sistema de plugins.
\end{itemize}

\subsection{Usuarios Objetivo}

\begin{itemize}
\item Productores de música chiptune.
\item Desarrolladores indie.
\item Estudiantes de ingeniería.
\end{itemize}

%====================================================
\section{Requisitos Funcionales}

%------------------------------------------------
\subsection{RF-01: Gestión de Proyectos}

\begin{itemize}
\item Crear proyecto nuevo.
\item Guardar proyecto.
\item Cargar proyecto.
\item Configurar BPM.
\item Configurar compás.
\end{itemize}

%------------------------------------------------
\subsection{RF-02: Gestión de Pistas}

\begin{itemize}
\item Crear/eliminar pistas.
\item Configurar instrumento por pista.
\item Ajustar volumen y paneo.
\item Asignar efectos.
\end{itemize}

%------------------------------------------------
\subsection{RF-03: Step Sequencer}

\begin{itemize}
\item Crear patrones de 8, 16 y 32 pasos.
\item Activar/desactivar pasos.
\item Configurar longitud de nota.
\item Soportar loop.
\end{itemize}

%------------------------------------------------
\subsection{RF-04: Motor de Síntesis 8-bit}

\begin{itemize}
\item Generar onda cuadrada.
\item Generar onda triangular.
\item Generar ruido blanco.
\item Aplicar envelope ADSR simplificado.
\end{itemize}

%------------------------------------------------
\subsection{RF-05: Reproducción}

\begin{itemize}
\item Play.
\item Stop.
\item Loop continuo.
\item Sincronización por clock interno.
\end{itemize}

%------------------------------------------------
\subsection{RF-06: Sistema MIDI}

\begin{itemize}
\item Detectar dispositivos MIDI.
\item Recibir Note On/Off.
\item Recibir Control Change.
\item Recibir Pitch Bend.
\item Grabar MIDI en tiempo real.
\item Cuantizar entrada MIDI.
\item Sincronizar mediante MIDI Clock.
\item Exportar archivo .mid.
\end{itemize}

%------------------------------------------------
\subsection{RF-07: Sistema de Plugins}

\begin{itemize}
\item Cargar plugins dinámicamente.
\item Registrar plugins automáticamente.
\item Soportar:
\begin{itemize}
    \item Instrument Plugins
    \item Effect Plugins
    \item Utility Plugins
\end{itemize}
\item Insertar plugins en cadena de efectos.
\item Eliminar plugins.
\end{itemize}

%------------------------------------------------
\subsection{RF-08: Audio Routing}

\begin{itemize}
\item Enviar pista a Track Bus.
\item Enviar pistas al Master Bus.
\item Soportar cadena de efectos por pista.
\end{itemize}

%------------------------------------------------
\subsection{RF-09: Exportación}

\begin{itemize}
\item Exportar WAV 8-bit.
\item Exportar MIDI.
\end{itemize}

%------------------------------------------------
\subsection{RF-10: Undo/Redo}

\begin{itemize}
\item Deshacer acción.
\item Rehacer acción.
\item Historial mínimo de 50 acciones.
\end{itemize}

%====================================================
\section{Requisitos No Funcionales}

\subsection{RNF-01: Arquitectura}

\begin{itemize}
\item Implementación bajo Clean Architecture.
\item Aplicación de principios SOLID.
\item Separación estricta GUI ↔ Dominio.
\end{itemize}

\subsection{RNF-02: Rendimiento}

\begin{itemize}
\item Latencia menor a 50ms.
\item Procesamiento eficiente de buffers.
\end{itemize}

\subsection{RNF-03: Extensibilidad}

\begin{itemize}
\item Nuevos plugins sin modificar núcleo.
\item Nuevos osciladores extensibles.
\end{itemize}

\subsection{RNF-04: Fiabilidad}

\begin{itemize}
\item Manejo robusto de excepciones.
\item Fallos en plugins no detienen el sistema.
\end{itemize}

\subsection{RNF-05: Usabilidad}

\begin{itemize}
\item Interfaz clara.
\item Flujo intuitivo de secuenciación.
\end{itemize}

\subsection{RNF-06: Portabilidad}

\begin{itemize}
\item Compatible con Windows 10+.
\end{itemize}

\subsection{RNF-07: Testabilidad}

\begin{itemize}
\item Cobertura mínima 75\%.
\item Tests unitarios por capa.
\end{itemize}

%====================================================
\section{Restricciones}

\begin{itemize}
\item Desarrollo inicial en Python.
\item Uso de PyQt6 para GUI.
\item Uso de sounddevice / PortAudio para audio.
\item Uso de mido + rtmidi para MIDI.
\end{itemize}

%====================================================
\section{Suposiciones}

\begin{itemize}
\item Usuario posee dispositivo MIDI opcional.
\item Sistema ejecutado en hardware estándar.
\end{itemize}

%====================================================
\section{Criterios de Aceptación Globales}

\begin{itemize}
\item El sistema reproduce patrones sin glitches.
\item El sistema detecta dispositivos MIDI correctamente.
\item Plugins pueden cargarse sin reiniciar la aplicación.
\item Exportación WAV produce archivo reproducible.
\end{itemize}

%====================================================
\section{Futuras Extensiones}

\begin{itemize}
\item Soporte VST.
\item Render offline avanzado.
\item Automatización avanzada.
\item Sincronización externa avanzada.
\end{itemize}

%====================================================
\section{Conclusión}

Este documento define formalmente los requisitos del sistema JustBeat-DAW, garantizando una base sólida para su diseño, implementación y evolución futura como plataforma profesional de producción musical 8-bit.

\end{document}